\scanpage{31}

il s'agit de décrire un iso. can.

\[
P(\underline{\operatorname{End}}(E)/\mathcal A)\simeq\deleted{\unsure{Q_{\mathrm{red}}}}\inserted{C,\ \text{conique can. dans }P(\underline{\operatorname{End}}^+(\mathcal M))}
\]

(\uline{NB} en termes de $E$ cependant, non en termes de $\mathcal A\subset\mathcal M$ seulement, car $\deleted{\unsure{Q}}\inserted{C}$ ne sera en général \unsure{pas} isom à un $P$(fibré vect. de rang $2$) \ldots)

Or on a sur $\underline{\operatorname{End}}(E)$ \uline{deux} structures de $\mathcal A$-module, la gauche et la droite, échangées par l'antiinvolution canonique

\[
u\longrightarrow\operatorname{Tr}(u)-u\qquad(\operatorname{Tr}=\text{trace réduite})
\]

et on a donc égalité des modules déterminants de $\underline{\operatorname{End}}(\deleted{\unsure{A}}E)$ sur \deleted{\unsure{a}} $\mathcal A$, lesquels sont égaux can. à $\underline{\operatorname{End}}(E)/\mathcal A$

\[
\underline{\operatorname{End}}(E)/\mathcal A\simeq\det_{\mathcal A}\underline{\operatorname{End}}(E)
\]

Mais l'iso

\[
\underline{\operatorname{End}}(E)\simeq E\otimes E^\vee
\]

transforme la structure de $\mathcal A$-module gauche en celle définie par l'opération $u_E\otimes\mathrm{id}_{E^\vee}$ de $u$, et la droite en $\mathrm{id}_E\otimes{}^tu_E$, \unsure{On|Donc} a \unsure{donc|ainsi|dans} \unsure{des iso|des isos|le cas} de $A$-modules

\[
\det_{\mathcal A}\operatorname{End}(E)\simeq E^{\otimes_A2}\otimes\underbrace{\bigl(\det_{\unsure{k|h}}E^\vee\bigr)}_{L^\vee}^{2}
\]

et de $=$

\[
\det_{\mathcal A}\operatorname{End}(E)\simeq\underbrace{\det_k(E)^{2}}_{L}\otimes {E^\vee}^{\otimes_A2},
\]

d'où un iso de $A$-modules

\[
\left[\operatorname{End}(E)/\mathcal A\simeq E^{\otimes_A2}\otimes_k L^\vee\simeq {E^\vee}^{\otimes_A2}\otimes_k L\right.
\]
% archive sheet number: 30
